\section{Zadanie 13 (Dominik Korsa)}
Dana jest moneta asymetryczna, która z nieznanym Ci prawdopodobieństwem
\(p \in (0, 1)\) zwraca orła i z prawdopodobieństwem \(1 - p\) zwraca reszkę.
Zaproponuj algorytm, który będzie symulował zachowanie monety symetrycznej,
to znaczy z prawdopodobieństwem \(\frac{1}{2}\) zwróci orła bądź reszkę.
Jedyne co możesz robić to rzucać monetą asymetryczną i na bazie obserwowanych
rezultatów kontynuować rzucanie bądź przerwać i zwrócić wynik.
Jaka jest oczekiwana liczba rzutów monetą asymetryczną przed zwróceniem wyniku?
Wynik podaj jako funkcję \(p\).

\subsection*{Algorytm}
\begin{itemize}
  \item Rzuć monetą dwa razy, wynik \(i\)-tego rzutu to \(e_i \in \{\texttt{OO}, \texttt{OR}, \texttt{RO}, \texttt{RR}\}\).
  \item Jeżeli \(e_i = \texttt{OR}\) to zwróć \texttt{R}.
  \item Jeżeli \(e_i = \texttt{RO}\) to zwróć \texttt{O}.
  \item Jeżeli \(e_i = \texttt{OO}\) lub \(e_i = \texttt{RR}\) to powtórz algorytm.
\end{itemize}

\subsection*{Model probabilistyczny}
\(\Omega = \{\texttt{OO}, \texttt{OR}, \texttt{RO}, \texttt{RR}\}^\mathbb{N}\)
-- zdarzeniami w tym modelu są ciągi par rzutów.\\
\(\mathcal{F} = \mathcal{P}(\Omega)\)

\(\{ A_i \}_{i \in \mathbb{N}}\) to rodzina zdarzeń niezależnych, gdzie \(A_i\) to zdarzenie,
że \(i\)-ta para rzutów daje \texttt{OR} lub \texttt{RO}, czyli kończy algorytm:
\[
  A_i = \Big\{ e \in \Omega \mid e_i \in \{\texttt{OR}, \texttt{RO}\} \Big\}
\]
\begin{align*}
  \forall_{i \in \mathbb{N}}\ &P(e_i = \texttt{OO}) = p^2\\
  \forall_{i \in \mathbb{N}}\ &P(e_i = \texttt{RR}) = (1-p)^2\\
  \forall_{i \in \mathbb{N}}\ &P(A_i) = p(1-p) + (1-p)p = 2p(1-p)
\end{align*}
Takie zdarzenia są niezależne.
Oba zdarzenia kończące algorytm są równoprawdopodobne, więc algorytm zwróci
\texttt{O} i \texttt{R} z takim samym prawdopodobieństwem.

Wprowadźmy zmienną losową \(X \colon \Omega \to \mathbb{N}\) oznaczającą liczbę
par rzutów wykonanych przez algorytm. \(P(X = 0) = 0\).
Dla \(n \in \mathbb{N}_+\) mamy:
\begin{align*}
P(X \geq n) &= P\big(
  \underbrace{(\Omega \setminus A_0) \cap (\Omega \setminus A_1) \cap \cdots \cap (\Omega \setminus A_{n-1})}
  _\text{pierwsze \(n-1\) rzutów nie zakończyło algorytmu}
\big)\\
&= \big(1-P(A_0)\big)\big(1-P(A_1)\big)\cdots\big(1-P(A_{n-1})\big) &&\text{(z niezależności)}\\
&= \big(1-2p(1-p)\big)^{n-1}
\end{align*}

\(X\) przyjmuje jedynie wartości w liczbach naturalnych więc z lematu P\&C
z Sufferingu mamy:
\begin{align*}
\mathbb{E}[X]
&= \sum_{n=1}^{\infty} P(X \geq n)
= \sum_{n=1}^{\infty} \big(1-2p(1-p)\big)^{n-1}
= \sum_{n=0}^{\infty} \big(1-2p(1-p)\big)^n\\
&= \frac{1}{2p(1-p)}
\end{align*}
\(\mathbb{E}[X] = \frac{1}{2p(1-p)}\) to oczekiwana liczba par rzutów,
więc oczekiwana liczba rzutów to:
\[
\boxed{\frac{1}{p(1-p)}}
\]
